Part~(i) is an exact criterion, but it is used as a decision rule, and a
decision rule is applied to estimates. Both $\varrho$ and $e_1/e_2$ must be
formed on data the decision does not see, so what governs the rule in
practice is not the criterion but the gap it has to resolve relative to the
noise in that estimation. The following makes this quantitative, and is
stated for any threshold comparison of this shape.

\begin{proposition}[margin law for a threshold rule]\label{prop:margin}
Let $D=\varrho-e_1/e_2$ be the signed gap of Proposition~\ref{prop:secmom}(i),
so that mixing the weaker member strictly helps exactly when $D<0$, and
suppose $D\neq0$. Let $\widehat D$ be an estimate of $D$ formed on a split
independent of the one on which $D$ is evaluated, and suppose the estimation
error $\widehat D-D$ is subgaussian with variance proxy $\sigma^2$, that is
$\mathbb E\exp\!\big(\lambda(\widehat D-D)\big)\le\exp(\lambda^2\sigma^2/2)$
for all $\lambda\in\mathbb R$, conditionally on $D$. Then, writing
$\Delta=|D|$ for the margin,
\begin{enumerate}
\item[(i)] on $\{D\neq0\}$, the decision taken from $\widehat D$ differs from
the correct one with probability at most
$\exp(-\Delta^{2}/2\sigma^{2})$;
\item[(ii)] for every $\tau\ge0$, and with no condition on $D$,
\[
\mathbb P\big(\operatorname{sign}\widehat D\neq\operatorname{sign}D,\ \
|\widehat D|\ge\tau\big)\;\le\;\exp\!\left(-\frac{\tau^{2}}{2\sigma^{2}}\right).
\]
\end{enumerate}
\end{proposition}

\begin{proof}
Suppose first $D<0$, so $\Delta=-D$ and the correct decision is to mix. The
decision taken from $\widehat D$ differs only if $\widehat D\ge0$, which
forces $\widehat D-D\ge-D=\Delta$. The subgaussian hypothesis and the
Chernoff bound give, for every $\lambda>0$,
$\mathbb P(\widehat D-D\ge\Delta)\le\exp(\lambda^2\sigma^2/2-\lambda\Delta)$,
and optimizing at $\lambda=\Delta/\sigma^2$ yields
$\exp(-\Delta^2/2\sigma^2)$. If instead $D>0$ then $\Delta=D$, the decision
differs only if $\widehat D\le0$, which forces $\widehat D-D\le-\Delta$, and
the same argument applied to $-(\widehat D-D)$, which is subgaussian with the
same proxy, gives the identical bound. The two cases are exhaustive because
$D\neq0$, which proves (i).

For (ii), condition on $D$ again. If $D<0$ the error event is
$\{\widehat D\ge0\}$, so on the error event intersected with
$\{|\widehat D|\ge\tau\}$ we have $\widehat D\ge\tau$ and therefore
$\widehat D-D\ge\tau-D\ge\tau$; the one-sided estimate above applies with
$\Delta$ replaced by $\tau$. If $D\ge0$ the error event is
$\{\widehat D<0\}$, so on it with $|\widehat D|\ge\tau$ we have
$\widehat D\le-\tau$ and therefore $D-\widehat D\ge\tau+D\ge\tau$, and the
one-sided estimate in the other direction applies. In both cases the bound is
$\exp(-\tau^2/2\sigma^2)$, which is free of $D$, so taking the expectation
over $D$ removes the conditioning.
\end{proof}

Three remarks on how the statement is used in
Section~\ref{sec:secmom-margin}. Part~(i) is conditional on the true margin
$\Delta$, which is not available when the decision is taken; what is
available is $|\widehat D|$, and part~(ii) is the statement about that
quantity. It is a joint bound: it limits the probability that a decision is
both wrong and taken at an observed margin of at least $\tau$, and it makes
that event rare among all decisions once $\tau$ exceeds a few multiples of
$\sigma$. It is not a bound on the error rate among the decisions whose
observed margin exceeds $\tau$ --- that conditional probability is the joint
one divided by $\mathbb P(|\widehat D|\ge\tau)$, on which the proposition
says nothing --- so the accuracy within a margin bin is an empirical
calibration of the rule and not a guarantee, and that is how
Section~\ref{sec:secmom-margin} reports it. Both bounds are vacuous for
margins of order $\sigma$ and fall off quickly beyond, so a rule of this
shape has a band around its own threshold inside which it carries no
guarantee, and the band is set by $\sigma$ rather than chosen. Second,
$\sigma$ is estimable without knowing $D$: the differences between
$\widehat D$ and the value of $D$ realized on a held-out block, taken across
many pairs, have exactly this scale. Third, $|\widehat D|$ is observable at
decision time, so a practitioner can report the observed margin next to the
verdict; part~(ii) says that verdicts reported at a large margin are seldom
wrong, and the scoreboard of Section~\ref{sec:secmom-margin} says how the
realized accuracy varies with that margin.


A second consequence of the same decomposition explains the final OCO-2
surrogate. Evaluation metrics there are diagonal in the output coordinates
(the radiance metric is the weighting $s_z$), and diagonal metrics decouple.

\begin{proposition}[per-coordinate combination under diagonal metrics]
\label{prop:percoord}
For weights $v=(v_{mj})$ that may depend on the output coordinate, let
$f_v(u)_j=\sum_m v_{mj}f_m(u)_j$. For any diagonal metric with positive
weights $w$,
\[
\mathbb E\big\|\mathrm{diag}(w)\big(f_v(u)-G(u)\big)\big\|_2^2
=\sum_j w_j^2\;\mathbb E\big(f_v(u)_j-G(u)_j\big)^2 ,
\]
so the minimizing weights solve $q$ separate problems, one per output
coordinate, none of which involves $w$. The per-coordinate optimum is the
same for every diagonal metric, and it weakly dominates every combination
whose weights are shared across coordinates, for all such metrics
simultaneously.
\end{proposition}

The practical content: when a problem is scored in several diagonal metrics
at once, the combination stage does not need to know which one matters. Fit
the best combination coordinate by coordinate on validation and the result
serves all of them; only the members need metric-aware training, which is
where the weighted mean of Section~\ref{sec:oco2} enters.

The proposition as stated covers quadratic metrics with fixed diagonal
weights, while the experiments report relative errors, which weight each
sample by its output norm. The decoupling survives that weighting.

\begin{corollary}[weighted metrics and the reported error]\label{cor:percoordw}
Let $a:\mathcal U\to[0,\infty)$ be measurable with
$\mathbb E[a(u)\,\|f_v(u)-G(u)\|_2^2]<\infty$. Then, for every diagonal $w$,
\[
\mathbb E\Big[a(u)\,\big\|\mathrm{diag}(w)\big(f_v(u)-G(u)\big)\big\|_2^2\Big]
=\sum_j w_j^2\;\mathbb E\Big[a(u)\,\big(f_v(u)_j-G(u)_j\big)^2\Big],
\]
so the minimization still splits into $q$ per-coordinate problems, none
involving $w$. Taking $a(u)=\|G(u)\|_2^{-2}$ and $w=\mathbf 1$ makes the left
side the mean squared relative error: the combination fitted coordinate by
coordinate with per-sample weights $a$ on validation is simultaneously
optimal, within its class, for the squared relative error and for every
diagonal reweighting of it. The reported metric, $\mathbb E\,\|\rho_{f_v}\|_2$
with $\rho$ the normalized residual of Section~\ref{sec:theory-secmom}, is
controlled through $\mathbb E\|\rho\|_2\le(\mathbb E\|\rho\|_2^2)^{1/2}$; the
gap between the two sides is the dispersion factor whose measured value stays
near $0.93$ across the pipelines of this paper.
\end{corollary}

The experiments run the combination both ways, unweighted and with the
per-sample weights of the corollary, and report both.

Proposition~\ref{prop:percoord} and its corollary describe the population
optimum; the pipeline computes an empirical selection, coordinate by
coordinate, on the validation split. The finite-sample cost of that step is
priced by the same argument that priced the stacks.

\begin{proposition}[empirical per-coordinate selection]\label{prop:selectcoord}
Let $f_1,\dots,f_M$ be fixed maps, $a:\mathcal U\to[0,A]$ a bounded
per-sample weight, and suppose $a(u)\,(f_m(u)_j-G(u)_j)^2\le L$ for
$\mu$-a.e.\ $u$, every member $m$ and coordinate $j$. On a validation sample
of size $m_v$ independent of the members, let $\widehat m(j)$ minimize the
empirical weighted risk of coordinate $j$ over the $M$ members, and let the
combination take coordinate $j$ from member $\widehat m(j)$. Then, with
probability at least $1-\delta$, simultaneously for every coordinate,
\[
R_j\big(\widehat m(j)\big)\;\le\;\min_{m\le M}R_j(m)\;+\;
2L\sqrt{\frac{\log(2Mq/\delta)}{2m_v}},
\qquad R_j(m)=\mathbb E\big[a(u)\big(f_m(u)_j-G(u)_j\big)^2\big],
\]
and consequently, for every diagonal metric $w$ at once,
$\sum_j w_j^2 R_j(\widehat m(j))\le\sum_j w_j^2\min_m R_j(m)
+2L\big(\sum_j w_j^2\big)\sqrt{\log(2Mq/\delta)/(2m_v)}$.
\end{proposition}

On the OCO-2 bands the deployed combination is exactly this estimator
($M=6$ members, $q=40$ coordinates, $m_v=2000$), in the unweighted form and,
for the squared relative error, with $a(u)=\|G(u)\|_2^{-2}$; the proposition
prices both at the same rate, and the selection cost it bounds is what the
seed experiments measure directly.

The corollary is the quantity this paper keeps measuring, at two different
levels. Across architecture families on the structural-mechanics benchmark
the pairwise correlations exceed $0.86$ at member errors near $4.7\%$, which
places the floor at $\ensFloor\pm\ensFloorSd\%$ in the root-mean-square
metric it is stated in; Section~\ref{sec:diversity} makes the comparison
with the deployed pipeline in that metric rather than across the two.
Across random seeds of one architecture on the OCO-2 bands the correlations
are $0.78$, $0.97$ and $0.96$ at member errors of $4.4\%$, $16.4\%$ and
$8.2\%$; measured in the squared metric of the proposition itself, over all
$45$ seed pairs per band, the floors are $4.360\%$, $17.970\%$ and $9.248\%$
and the realized ten-seed ensembles sit at $4.419\%$, $17.997\%$ and
$9.267\%$ (Section~\ref{sec:oco2}). When an ensemble is at its floor the
corollary says where improvement cannot come from. The reported tables use
the mean rather than the root mean square of $\|\rho_f\|_2$, and the two are
not interchangeable: Remark~\ref{rem:metric} gives the relation, the ratio
runs from about $0.89$ to $0.94$ across the objects here, and every
comparison in this paper between a quantity read off $S$ and a table entry is
made after converting one to the other. The proposition frames the diversity
experiment of Section~\ref{sec:diversity}: what a new member is worth is
legible in the row it adds to $S$, before any weights are fitted. On this
benchmark the rows all look alike, pairwise correlations of
\corrNetLo{}--\corrNetHi{} among trained networks of three architecture
families and two losses, \corrKerLo{}--\corrKerHi{} against the kernel
predictor, so the floor sits just below the best single member, and the
measured stack lands on it. Part (i) is the same
computation specialized to two members; it correctly predicts, from measured
$(e_1,e_2,\varrho)$ alone, which members receive weight in the fitted stack
and which are dropped. Section~\ref{sec:oco2} applies the same test on a
problem where the answer comes out the other way: there the kernel's error is
inflated by a factor the next subsection quantifies, the threshold $e_1/e_2$
collapses, and the kernel is dropped even at a residual correlation of $0.14$.


Suppose the operator factors as $G(u)=g(Au)$ with $A\in\mathbb R^{d\times d}$
nonsingular and $g$ of unit $H^{s}$ norm, and that $\mu$ has a density bounded
above and below on a compact domain. Order the singular values
$\sigma_1\ge\dots\ge\sigma_d$ of $A$; they are the rates at which $G$ varies
along the input directions. Suppose $A$ has \emph{approximate rank} $r$ at the
sample size, meaning the trailing singular values fall below the achievable
resolution, $\sigma_{r+1}\lesssim n^{-1/r}$, so the image $A\,\mathcal U$ is
effectively $r$-dimensional at scale $n^{-1/r}$. The idealization beneath the
display also takes the leading-$r$ projection of the image design to keep a
density bounded above and below, the trailing directions to be exact zeros at
the working resolution, and the interpolation limit $\lambda\to0$;
Supplement~\ref{app:proofs} records the argument at that level, as a sketch
rather than a complete proof. The statement is accordingly a heuristic rate
calculation, and it is labelled informal: nothing else in the paper rests on
it, and its role is to say what kind of gain a metric can and cannot produce.

\begin{proposition}[isotropic and adapted kernel rates, informal]\label{prop:aniso}
Let $\widehat G^{\mathrm{iso}}_n$ be kernel ridge regression with a Mat\'ern
kernel of smoothness $s$ and a validation-optimal single length scale, and
$\widehat G^{\mathrm{ard}}_n$ the same construction in the metric
$u\mapsto Au$. Up to constants depending on $(\sigma_i)$, $g$, $s$ and the
density bounds, and up to logarithmic factors in $n$, with high probability
over $n$ i.i.d.\ samples,
\[
\mathbb E\big\|G-\widehat G^{\mathrm{iso}}_n\big\|\;\lesssim\;\sigma_1^{\,s}\,n^{-s/d},
\qquad
\mathbb E\big\|G-\widehat G^{\mathrm{ard}}_n\big\|\;\lesssim\;n^{-s/r},
\]
so the ratio of the two bounds is of order $\sigma_1^{\,s}\,n^{\,s(1/r-1/d)}$.
When $A$ is close to full rank ($r\approx d$) the two rates coincide and only
the constant $\sigma_1^{\,s}$ separates them.
\end{proposition}

The calculation (Supplement~\ref{app:proofs}) is the scattered-data estimate
$\|G-\widehat G_n\|\lesssim h_X^{\,s}\|G\|_{H^s}$ applied to the two designs:
one length scale must fill all $d$ directions, so $h_X\asymp n^{-1/d}$ and the
norm carries $\sigma_1^{\,s}$; the adapted metric, when $A$ has a spectral gap,
confines the design to $r$ resolved directions and improves the fill distance.
Both displays are upper bounds under the stated idealization --- the exact
metric $A$, a fixed $g$, and the length-scale selection folded into the
logarithmic factors --- and we claim no matching lower bound; the heuristic
is used here to calibrate what a metric can and cannot buy, not as a rate
theorem.
It isolates the \emph{metric} part of a network's advantage, the
part an anisotropic kernel would recover, and it says this part is a rate gain
only under an approximate-rank gap; without one it is the constant
$\sigma_1^{\,s}$ alone. The rest is \emph{representational}: a stationary
kernel of fixed smoothness reaches only its native space, at any metric, while
the network adapts its features to the map. The two parts are separated
experimentally by handing the kernel the anisotropic metric and seeing how
much of the gap closes. On the OCO-2 bands of Section~\ref{sec:oco2} the metric
closes a factor of $1.1$ to $1.5$ out of a tenfold gap
(Section~\ref{sec:aniso-measured}); the remainder is representational, and its
direct evidence is that the same kernel machinery applied to the network's
learned features recovers the network's accuracy and slightly exceeds it.

The rate half of the heuristic is tested rather than assumed, and the test is
the more interesting for not going our way. Both kernel rows are refit at
nested training sizes over a sixteenfold range, at ten seeds per band. The
sensitivity map's approximate rank is $0.54$ to $0.60$ of the input dimension,
a gap wide enough that the calculation would permit a rate gain; the measured
difference between the two empirical exponents is nevertheless small on all
three bands and consistent with zero on one. The displays are upper bounds
with no matching lower bound, so this does not contradict them --- the adapted
rate is free to beat its own bound --- but it does mean that on these problems
the rank gap is not converted into a rate, and a reader should take from
Proposition~\ref{prop:aniso} only the part the measurement supports: an
anisotropic metric is worth a constant. The controlled rank-gap target of the
verification suite, where the construction is exactly the idealization the
calculation assumes, does show the adapted exponent improving; the distance
between that and the OCO-2 bands is the distance between the idealization and
a real sensitivity map.


